% !TeX program = xelatex
\documentclass[11pt,a4paper]{article}
\usepackage{fontspec}
\setmainfont{DejaVu Serif}
\setsansfont{DejaVu Sans}
\setmonofont{DejaVu Sans Mono}
\renewcommand{\contentsname}{Spis treści}
\renewcommand{\tablename}{Tabela}
\hyphenpenalty=3000 \tolerance=2000 \emergencystretch=2em
\usepackage[margin=2.1cm]{geometry}
\usepackage{amsmath}
\usepackage{xcolor}
\definecolor{granat}{RGB}{12,35,88}
\usepackage{booktabs}
\usepackage{tabularx}
\usepackage{longtable}
\usepackage{array}
\usepackage{enumitem}
\usepackage[colorlinks=true,linkcolor=granat,urlcolor=granat]{hyperref}
\usepackage{titlesec}
\usepackage{parskip}
\titleformat{\section}{\large\bfseries\sffamily\color{granat}}{\thesection.}{0.6em}{}
\titleformat{\subsection}{\normalsize\bfseries\sffamily\color{granat}}{\thesubsection.}{0.6em}{}
\setlist{nosep,leftmargin=1.4em}
\newcolumntype{L}{>{\raggedright\arraybackslash}X}
% skróty typograficzne
\newcommand{\ui}[1]{\textbf{#1}}                    % element interfejsu
\newcommand{\szukaj}[1]{\mbox{\ui{Szukaj polecenia} $\to$ \emph{#1}}} % fallback
\newcommand{\krok}[1]{\par\medskip\noindent\textbf{\color{granat}Krok #1.}\space}

\title{\textbf{\color{granat}AMPERE POINT --- custom device}\\[4pt]
\large Przewodnik EPLAN krok po kroku:\\ odtworzenie schematu elektrycznego v3 w EPLAN Education 2026\\[2pt]
\normalsize wersja 1 --- dla osoby bez doświadczenia z EPLAN}
\author{Opracowanie wewnętrzne AMPERE POINT}
\date{6 lipca 2026}

\begin{document}
\maketitle
\vspace{-1.6em}

\noindent\textbf{Cel:} przenieść do EPLAN trzy arkusze schematu elektrycznego custom device
(wersja 3: \texttt{AMPERE\_POINT\_custom\_device\_schemat\_elektryczny\_v3\_EPLAN.pdf} --- miej ten
PDF otwarty obok jako wzór rysunkowy). \textbf{Stan wyjściowy:} dokładnie ten z Twojego ekranu ---
EPLAN Education 2026 (interfejs polski), otwarty projekt \texttt{AMPERE\_POINT\_custom\_device}
(AP-CD-001) ze stroną tytułową \texttt{=CA1+EAA/1}.

\textbf{Jak pisany jest ten przewodnik (ważne):} nazwy kart wstążki (\ui{Plik},
\ui{Narzędzia główne}, \ui{Wstaw}, \ui{Edytuj}, \ui{Widok}, \ui{Narzędzia}\ldots) oraz elementy
okna (\ui{Centrum wprowadzania}, nawigator \ui{Strony}, pole \ui{Szukaj polecenia}, pasek stanu
z siatką) są zgodne z Twoim zrzutem ekranu; ścieżki tworzenia stron i typy stron sprawdzono
w oficjalnej pomocy EPLAN. Ponieważ drobne nazwy przycisków bywają różne między wydaniami,
\emph{każda} instrukcja ma niezawodny wariant zapasowy: pole \ui{Szukaj polecenia} (lupka na
górnym pasku okna) --- wpisujesz tam podaną nazwę i EPLAN sam pokazuje oraz uruchamia właściwą
funkcję. Jeśli więc jakiegoś przycisku nie widzisz na wstążce --- nie szukaj po omacku, użyj
\ui{Szukaj polecenia}.

\tableofcontents

% =====================================================================
\section{Trzy zasady, zanim klikniesz cokolwiek}

\begin{enumerate}
\item \textbf{EPLAN nie ma przycisku ,,Zapisz''.} Każda zmiana na stronie zapisywana jest do bazy
projektu od razu. Nie szukaj Ctrl+S; pomyłki wycofujesz przez \ui{Cofnij} (Ctrl+Z, strzałki na
pasku szybkiego dostępu w lewym górnym rogu).
\item \textbf{Wszystko wstawia się ,,na kursor''.} Po wybraniu symbolu/polecenia symbol wisi na
kursorze; klik lewym = wstawienie (można wstawiać wielokrotnie), \textbf{Esc} = koniec.
W trakcie wstawiania klawisz \textbf{Tab} obraca/zmienia wariant symbolu.
\item \textbf{Dwuklik = właściwości.} Każdy wstawiony element (symbol, tekst, skrzynka) otwiera
dwuklikiem okno właściwości --- tam zmieniasz oznaczenie (DT), numery przyłączy, teksty.
Zaznaczenie + \textbf{Delete} usuwa.
\end{enumerate}

% =====================================================================
\section{Orientacja w oknie (mapa Twojego ekranu)}

\begin{itemize}
\item \textbf{Wstążka} u góry --- karty \ui{Plik}, \ui{Narzędzia główne}, \ui{Wstaw},
\ui{Edytuj}, \ui{Widok}, \ui{Połączenia}, \ui{Narzędzia}, \ui{Dane zasadnicze}. Większość rysowania
robi się z karty \ui{Wstaw} albo z Centrum wprowadzania.
\item \textbf{Nawigator Strony} (lewy panel, zakładka z nazwą projektu) --- drzewo stron:
\texttt{CA1 $\to$ EAA $\to$ 1 Strona tytułowa}. Dwuklik na stronie otwiera ją w edytorze.
\item \textbf{Edytor graficzny} (środek) --- tu rysujesz; na dole pasek stanu ze współrzędnymi
X/Y, \ui{Siatka C: 4,00 mm} i przełącznikami przyciągania. Zostaw siatkę 4 mm --- to domyślna
siatka stron schematów i wszystkie odległości w tym przewodniku jej odpowiadają.
\item \textbf{Centrum wprowadzania} (prawy panel) --- najważniejsze narzędzie tego przewodnika:
pole \ui{Szukaj} + sekcje \ui{Symbole}, \ui{Urządzenia}, \ui{Makra okna/symbolu},
\ui{Ostatnio używane}, \ui{Ulubione}. Wpisujesz polskie hasło (np. \emph{bezpiecznik}),
klikasz wynik i wstawiasz na stronę.
\item \textbf{Podgląd graficzny} (lewy dół) --- miniatura aktualnej strony.
\item Zamknięty panel przywrócisz wpisując jego nazwę w \ui{Szukaj polecenia}
(np. \emph{Strony}, \emph{Centrum wprowadzania}, \emph{Podgląd graficzny}).
\end{itemize}

Nawigacja w edytorze: kółko myszy = zoom; przytrzymanie kółka (środkowego przycisku) =
przesuwanie widoku; klawisz \textbf{3} na klawiaturze = dopasuj całą stronę do okna
(jeśli nie zadziała: \szukaj{Cała strona}).

% =====================================================================
\section{Utworzenie stron E1--E3}
\label{sec:strony}

Schemat v3 ma trzy arkusze robocze. W strukturze projektu będą to strony
\texttt{=CA1+EAA/2}, \texttt{/3} i \texttt{/4} (strona \texttt{/1} to istniejąca tytułowa).
Oznaczenia ,,E1/E2/E3'' zapiszemy w opisie strony.

\krok{1} W nawigatorze \ui{Strony} kliknij \textbf{prawym} przyciskiem na węźle \ui{EAA}
i wybierz \ui{Nowy\ldots} (to samo daje karta \ui{Narzędzia główne} $\to$ grupa \ui{Strona}
$\to$ \ui{Nowy}). Otworzy się okno \ui{Nowa strona}.

\krok{2} Wypełnij:
\begin{itemize}
\item \ui{Pełna nazwa strony}: \texttt{=CA1+EAA/2} (przycisk \ui{\ldots} obok pola otwiera
pomocnika struktury, ale można wpisać ręcznie);
\item \ui{Typ strony}: \textbf{Schemat wielobiegunowy (I)} --- wybierz z listy rozwijanej;
litera (I) oznacza stronę interaktywną, rysowaną ręcznie;
\item \ui{Opis strony}: \texttt{E1 --- tor mocy 3F};
\item resztę zostaw domyślną (siatka 4 mm jest domyślna dla tego typu strony) i kliknij \ui{OK}.
\end{itemize}

\krok{3} Powtórz dla dwóch kolejnych stron:
\texttt{/3} z opisem \texttt{E2 --- łańcuch bezpieczeństwa 24 V} oraz
\texttt{/4} z opisem \texttt{E3 --- obciążenie i sterowanie}.

\krok{4} Sprawdź w drzewie: pod \ui{EAA} są teraz strony 1--4. Dwuklik na \texttt{/2} otwiera
pustą stronę z ramką: kolumny 1--10 u góry/dołu i wiersze A--F po bokach pochodzą z formularza
ramki projektu. \textbf{Tak czytamy odsyłacze:} zapis \texttt{/E2.3} z PDF-a v3 oznacza
,,strona E2 (czyli /3), kolumna 3''.

\emph{Uwaga:} jeśli w Twoim szablonie okno \ui{Nowa strona} pokazuje od razu inne pola
(np. skalę) --- nie zmieniaj ich. Liczy się typ strony, nazwa i opis.

% =====================================================================
\section{Elementarz rysowania --- sześć czynności, z których składa się wszystko}
\label{sec:elementarz}

Cały schemat powstaje z powtarzania sześciu czynności. Dalsze rozdziały odsyłają do nich
numerami (E.1--E.6).

\subsection*{E.1 --- Wstawienie symbolu z Centrum wprowadzania}
\begin{enumerate}
\item W \ui{Centrum wprowadzania} (prawy panel) kliknij \ui{Symbole} (albo od razu wpisz hasło
w pole \ui{Szukaj}, np. \emph{bezpiecznik}).
\item Kliknij znaleziony symbol --- zawiśnie na kursorze. \textbf{Tab} przełącza warianty
(obrót, odbicie), kółko przybliża widok.
\item Kliknij w miejscu docelowym na stronie. Otworzy się okno \ui{Właściwości}:
  \begin{itemize}
  \item \ui{Widoczne oznaczenie DT} --- wpisz oznaczenie z wzoru, np. \texttt{-F1}
  (z myślnikiem na początku, jak w PDF);
  \item zakładka/pola przyłączy --- numery zacisków, np. \texttt{1} i \texttt{2}
  (u nas: 13/14 dla zwiernych pomocniczych, 11/12 i 21/22 dla rozwiernych, A1/A2 dla cewek);
  \item \ui{OK}.
  \end{itemize}
\item Symbol nadal wisi na kursorze (można wstawić następny egzemplarz); \textbf{Esc} kończy.
\end{enumerate}
Alternatywa: karta \ui{Wstaw} $\to$ \ui{Symbol} otwiera pełną przeglądarkę bibliotek
(IEC\_symbol); Centrum wprowadzania jest po prostu szybsze.

\subsection*{E.2 --- Połączenia (autoconnecting)}
EPLAN \textbf{sam} rysuje połączenie, gdy dwa punkty przyłączeń (czerwone strzałki na końcach
symboli) znajdą się dokładnie w jednej pionowej lub poziomej linii siatki. Nie rysuje się więc
,,kabli'' ręcznie --- ustawia się symbole w osi. Jeśli połączenie nie powstało: element leży poza
siatką (przesuń go; ma się ,,przyciągać'' co 4 mm) albo przyłącza nie są naprzeciw siebie.

\subsection*{E.3 --- Zagięcia i rozgałęzienia (symbole połączeń)}
Do prowadzenia linii pod kątem i do odczepów służą \emph{symbole połączeń}:
\ui{Narożnik} (zagięcie 90$^\circ$) i \ui{Trójnik} (rozgałęzienie --- kropka węzła).
Znajdziesz je na karcie \ui{Wstaw}; najpewniej przez \szukaj{Narożnik} oraz \szukaj{Trójnik}.
Wstawiasz je \emph{na} istniejącej linii/osi połączenia; Tab obraca kierunek odejścia.
Skrzyżowanie dwóch linii bez trójnika = brak połączenia (tak mają się krzyżować np. odczepy
-X3 z pionami L2/L3 na E1).

\subsection*{E.4 --- Punkt przerwania (kontynuacja potencjału na innej stronie)}
,,Punkt przerwania'' to strzałka z nazwą (np. \texttt{+24V}): dwa punkty przerwania o tej samej
nazwie w projekcie łączą się logicznie i pokazują nawzajem swoje adresy (dokładnie jak
odsyłacze \texttt{/E3.5} w PDF). Wstawianie: \szukaj{Punkt przerwania}, klik na końcu linii,
w właściwościach wpisz \ui{Widoczne oznaczenie DT} = nazwę potencjału (np. \texttt{+24V},
\texttt{0V}, \texttt{L1\_MB}). Używamy ich do: szyn 24 V (E2$\leftrightarrow$E3), zasilania
modułu B (E1$\to$E3) i odczepu generatora (E1$\to$E3).

\subsection*{E.5 --- Czarna skrzynka (urządzenie bez gotowego symbolu)}
Bloki takie jak -A1 (płyta ESP32), -A3 (generator), -G1 (zasilacz), -B1 (czujnik CDSR),
-A2 (watchdog), -P1 (HMI) rysuje się jako \emph{czarną skrzynkę}: prostokąt reprezentujący
urządzenie. Wstawianie: \szukaj{Czarna skrzynka} (na karcie \ui{Wstaw});
kliknij pierwszy narożnik, potem przeciwległy; w właściwościach nadaj DT (np. \texttt{-A1})
i ewentualnie tekst opisu. Jeżeli skrzynka ma mieć własne, podłączalne piny, dodaj na jej
krawędzi \emph{punkty przyłączenia urządzeń}: \szukaj{Punkt przyłączenia} --- każdy dostaje numer
(np. 1, 2\ldots). Dla naszych celów w większości wystarczą skrzynki z tekstami, jak w PDF.

\subsection*{E.6 --- Teksty, grafika, powielanie}
\begin{itemize}
\item \textbf{Tekst wolny} (uwagi, opisy przy skrzynkach): \ui{Wstaw} $\to$ \ui{Tekst}
(lub \szukaj{Tekst}); klik w miejscu docelowym, wpisz treść, OK.
\item \textbf{Prostokąt/linia jako grafika} (np. przerywana obwiednia ,,MODUŁ B''):
\ui{Wstaw} $\to$ grupa grafiki $\to$ \ui{Prostokąt} (\szukaj{Prostokąt}); po narysowaniu
dwuklik $\to$ w formatowaniu ustaw \ui{Rodzaj linii} przerywany. Grafika nie tworzy połączeń
elektrycznych --- można nią bezpiecznie obrysowywać.
\item \textbf{Powielanie} (klucz do 12 gałęzi load banku): zaznacz fragment (przeciągnij ramkę),
karta \ui{Edytuj} $\to$ \ui{Powiel} (widoczny na Twoim zrzucie), wskaż punkt odniesienia
i miejsca kopii; EPLAN zapyta/ponumeruje DT kolejnych kopii --- i tak sprawdź je dwuklikiem.
Zwykłe Ctrl+C/Ctrl+V też działa.
\end{itemize}

% =====================================================================
\section{Arkusz E1 (strona /2) --- tor mocy}
\label{sec:e1}

Otwórz stronę \texttt{/2} (dwuklik w nawigatorze). Wzór: strona 2/4 PDF-a v3.
Rysuj z grubsza w tych samych kolumnach ramki co wzór --- odsyłacze będą się zgadzać z PDF-em.

\textbf{Aparaty tego arkusza} (hasła do pola Szukaj w Centrum wprowadzania):

\begin{longtable}{@{}p{0.10\textwidth}p{0.26\textwidth}p{0.25\textwidth}p{0.31\textwidth}@{}}
\toprule
\textbf{DT} & \textbf{Co to} & \textbf{Hasło wyszukiwania} & \textbf{Uwagi}\\
\midrule
\endhead
-X1:1..7 & zaciski gniazda Type 2 & \emph{zacisk} & 7 szt. w poziomie, u góry\\
-X3:1..7 & zaciski odczepu PM701E & \emph{zacisk} & 7 szt. w pionie, po prawej\\
-K4 & zestyk zwierny w linii CP & \emph{zestyk zwierny} & DT \texttt{-K4}; cewka będzie na E2\\
-Q0 & rozłącznik 4-bieg. & \emph{rozłącznik} & wariant 3-bieg. + 1-bieg. z tym samym DT\\
-F1..-F3 & bezpieczniki gG 32 A & \emph{bezpiecznik} & po jednym na L1/L2/L3\\
-B1 & czujnik różnicowy CDSR & --- (czarna skrzynka, E.5) & obejmuje 4 piony, bez PE\\
-T1..-T3 & przekładniki prądowe & \emph{przekładnik prądowy} & na L1/L2/L3\\
-K1 & styki główne stycznika & \emph{zestyk mocy} lub \emph{stycznik} & 4 szt., DT \texttt{-K1}\\
-X2:1..5 & zaciski wyjścia CEE & \emph{zacisk} & na dole\\
\bottomrule
\end{longtable}

\krok{1} \textbf{Zaciski -X1.} Wstaw symbol zacisku (E.1) i kliknij 7 razy w jednej linii
poziomej u góry strony, w odstępach ok.\ 5--6 pól siatki. W właściwościach pierwszego wpisz
DT \texttt{-X1} i \ui{oznaczenie zacisku} \texttt{1}; kolejne EPLAN zwykle numeruje sam (sprawdź
dwuklikiem). Nad zaciskami dodaj teksty (E.6): L1, L2, L3, N, PE, CP, PP.

\krok{2} \textbf{Piony.} Zaciski mają przyłącza w osi pionowej --- poprowadź układ tak, aby
kolejne aparaty (kroki 4--8) stawały dokładnie pod zaciskami 1--4 (L1..N): wtedy połączenia
narysują się same (E.2). Pion PE (zacisk 5) zostaje pusty aż do -X2 --- po prostu seria
połączeń w dół.

\krok{3} \textbf{Odczep -X3 i zestyk -K4.} Na pionach L1, L2, L3, N (na wysokości tuż pod -X1)
wstaw \ui{Trójniki} (E.3) i wyprowadź poziome linie w prawo; zakończ je 7 zaciskami \texttt{-X3}
ustawionymi pionowo (numery 1--7). Linie CP i PP prowadź od zacisków 6 i 7 gniazda -X1;
\textbf{w linii CP wstaw zestyk zwierny} (E.1) z DT \texttt{-K4} (przyłącza 13/14). Odsyłacz przy
-K4 pojawi się automatycznie, gdy na E2 powstanie cewka -K4. Dodaj teksty: \emph{odczep Type 2
$\to$ PM701E (7 żył, $\le$1 m)}, \emph{SIGNAL OUTPUT (CP)}, \emph{mostek serwisowy CP na
listwie -X3}.

\krok{4} \textbf{-Q0.} Wyszukaj \emph{rozłącznik}; wybierz wariant 3-biegunowy i wstaw na pionach
L1--L3; obok, na pionie N, wstaw wariant 1-biegunowy i nadaj mu \textbf{ten sam DT} \texttt{-Q0}
--- EPLAN potraktuje oba jako jeden aparat (spyta lub scali w nawigatorze urządzeń; to normalne).
Numery przyłączy: 1/2, 3/4, 5/6, 7/8 jak w PDF.

\krok{5} \textbf{-F1..-F3.} Symbol \emph{bezpiecznik} na L1, L2, L3 (N bez zabezpieczenia ---
zostaje samo połączenie). Tekst obok: \emph{gG 32 A}.

\krok{6} \textbf{-B1 (CDSR).} Czarna skrzynka (E.5) narysowana \textbf{w poprzek} czterech pionów
L1--N (wąski, leżący prostokąt jak w PDF), DT \texttt{-B1}, tekst: \emph{czujnik różnicowy AC/DC
(CDSR 0.07-TP), okno L1+L2+L3+N} oraz strzałka/tekst \emph{SPI $\to$ /E3.3}. Piony przechodzą
przez skrzynkę --- to odwzorowanie ,,okna'' czujnika; nie potrzeba tu przyłączy elektrycznych.

\krok{7} \textbf{Odczep do generatora -A3 (poprawka v3).} Na pionie L1, \textbf{poniżej} skrzynki
-B1 a nad -T1, wstaw \ui{Trójnik}, krótką linię w lewo i \ui{Punkt przerwania} (E.4) o nazwie
\texttt{L1\_A3}; jego para znajdzie się na E3 przy generatorze. Tekst: \emph{przez -K30}.

\krok{8} \textbf{-T1..-T3, potem -K1, potem -X2.} Przekładniki (hasło: \emph{przekładnik prądowy})
na L1--L3; niżej cztery zestyki mocy z DT \texttt{-K1} (przyłącza 1/2\ldots7/8); na dole 5 zacisków
\texttt{-X2} (L1, L2, L3, N, PE --- tu pion PE dołącza). Teksty: \emph{40 A / 1 V},
\emph{6 mm$^2$}, \emph{gniazdo CEE 32 A 5p $\to$ moduł B}, \emph{przewód 5$\times$6 mm$^2$
H07RN-F}. Pod -K1 automatycznie pojawi się odsyłacz do cewki (gdy narysujesz ją na E2).

\krok{9} \textbf{Zasilanie modułu B na E3.} Zamiast prowadzić moc ,,poza stronę'', wstaw przy
-X2 cztery \ui{Punkty przerwania}: \texttt{L1\_MB}, \texttt{L2\_MB}, \texttt{L3\_MB},
\texttt{N\_MB} --- ich pary będą początkiem szyn modułu B na E3 (rozdz.~\ref{sec:e3g}).

\krok{10} \textbf{Uwagi tekstowe} (E.6) po prawej: przepisz bloki \emph{Uwagi},
\emph{Ostrzeżenie (v3)} i \emph{-K4 (v3)} z PDF-a. Gotowe --- porównaj stronę z wzorem
(klawisz 3 = cała strona).

% =====================================================================
\section{Arkusz E2 (strona /3) --- łańcuch bezpieczeństwa 24 V}
\label{sec:e2g}

Wzór: strona 3/4 PDF-a. To schemat drabinkowy: dwie poziome szyny (+24 V u góry, 0 V u dołu)
i pionowe ścieżki między nimi.

\textbf{Aparaty:} -S0 (E-STOP, 2 zestyki rozwierne 11/12 i 21/22), -F11.1 i -F11.12 (termiki,
zestyki rozwierne z wyzwalaczem termicznym), -S1 (krańcówka, rozwierny), -K2/-K10/-RB (zestyki
zwierne 13/14), -A2 (czarna skrzynka + zestyk), -S2 (przycisk zwierny), cewki: -RB, -K1, -K4.

\krok{1} \textbf{Szyny.} Wstaw \ui{Punkt przerwania} (E.4) o nazwie \texttt{+24V} przy lewej
krawędzi obszaru roboczego i drugi \texttt{+24V} przy prawej; ustaw je dokładnie w jednej osi
poziomej --- EPLAN sam pociągnie między nimi linię (to nasza szyna). Tak samo zbuduj szynę
\texttt{0V} na dole. (Punkty o tych samych nazwach wykorzystamy też na E3 przy zasilaczu -G1 ---
wtedy przy szynach pojawią się odsyłacze, jak \texttt{z G1 /E3.5} w PDF.)

\krok{2} \textbf{Ścieżka 1 (kolumna 1): E-STOP i termiki.} Z szyny +24 V w dół (Trójnik, E.3):
zestyk \textbf{rozwierny} (hasło: \emph{zestyk rozwierny}; jeśli w bibliotece jest wariant
,,przycisk grzybkowy/awaryjny'' --- użyj go) z DT \texttt{-S0}, przyłącza 11/12. Pod nim dwa
zestyki rozwierne \texttt{-F11.1} i \texttt{-F11.12} (jeśli znajdziesz wariant \emph{z wyzwalaczem
termicznym} --- lepiej; jeśli nie, zwykłe rozwierne + tekst $\theta$ obok) i wielokropek tekstem
,,\vdots'' między nimi --- rysujemy pierwszy i ostatni z dwunastu, jak w PDF. Niżej krańcówka
\texttt{-S1} (rozwierny; wariant z rolką, jeśli dostępny). Obrysuj termiki przerywanym
prostokątem (E.6) i podpisz: \emph{termiki KSD301 sekcji grzałek, NC szeregowo; moduł B, złącze
-X5}. Ścieżkę zakończ \ui{Punktem przerwania} \texttt{LB1} (przejście do ścieżki 3 --- w PDF
odpowiada strzałce /E2.3).

\krok{3} \textbf{Ścieżka 3 (kolumna 3): zezwolenia i cewka -RB.} Zacznij od pary punktu
przerwania \texttt{LB1}, dalej w dół: zestyki \textbf{zwierne} \texttt{-K2} (13/14),
\texttt{-K10} (13/14), \texttt{-A2}. Poniżej gałąź równoległa: Trójnik, w lewej odnodze
\emph{przycisk zwierny} \texttt{-S2} (START), w prawej zestyk zwierny \texttt{-RB} 13/14
(samopodtrzymanie), odnogi połącz Trójnikiem z powrotem. Pod spodem \textbf{cewka} (hasło:
\emph{cewka}) z DT \texttt{-RB} (A1/A2) i dalej do szyny 0 V. Podpisy tekstem jak w PDF
(\emph{potwierdzenie wentylatorów}, \emph{zezwolenie MCU}, \emph{watchdog impulsowy},
\emph{START}, \emph{samopodtrzymanie}, \emph{przekaźnik bezpieczeństwa}).

\krok{4} \textbf{Ścieżka 5: cewka -K1.} Z szyny +24 V przez zestyk zwierny \texttt{-RB}
(przyłącza \textbf{23/24}) do \textbf{cewki} \texttt{-K1}, dalej do 0 V. Gdy tylko nadasz cewce
DT \texttt{-K1}, pod cewką pojawi się \emph{lustro styków} --- automatyczny spis styków -K1
z E1 z adresami (to samo, co ręcznie wypisane w PDF). Jeśli lustro się nie pokazuje, otwórz
właściwości cewki i poszukaj ustawienia \emph{lustra styków} (wyświetlanie na cewce);
w razie trudności: \szukaj{lustro styków}.

\krok{5} \textbf{Ścieżka 6 (nowa w v3): cewka -K4.} Z szyny +24 V w dół: zestyk
\textbf{rozwierny} z DT \texttt{-S0} i przyłączami \textbf{21/22} (drugi zestyk tego samego
E-STOPa --- ten sam DT!), pod nim \textbf{cewka} \texttt{-K4}, dalej do 0 V. EPLAN pokaże przy
obu zestykach -S0 wzajemne odsyłacze, a przy cewce -K4 --- odsyłacz do zestyku CP z E1.
Tekst: \emph{przerwanie CP (v3)}.

\krok{6} \textbf{-A2 jako urządzenie.} Watchdog to fizyczny moduł: obok ścieżki (jak w PDF na E3)
wystarczy zestyk z DT \texttt{-A2} tutaj i czarna skrzynka \texttt{-A2} na E3; oba elementy
z tym samym DT zwiążą się odsyłaczem.

\krok{7} \textbf{Adnotacje} (E.6): przepisz bloki \emph{Zasada}, \emph{-K4 (v3)},
\emph{Odczyty diagnostyczne}, \emph{Poza łańcuchem} z PDF-a. Porównaj całość ze wzorem.

% =====================================================================
\section{Arkusz E3 (strona /4) --- obciążenie i sterowanie}
\label{sec:e3g}

Wzór: strona 4/4 PDF-a. Najwięcej elementów, ale w 80\% to powielanie.

\krok{1} \textbf{Szyny modułu B.} U góry strony wstaw pary punktów przerwania z E1:
\texttt{L1\_MB}, \texttt{L2\_MB}, \texttt{L3\_MB} (trzy poziome szyny jedna pod drugą)
oraz niżej \texttt{N\_MB} (czwarta szyna). Każdą szynę zakończ po prawej drugim punktem tej
samej nazwy albo po prostu pozostaw linię do ostatniej gałęzi. Obrysuj całość przerywanym
prostokątem (E.6) i podpisz: \emph{MODUŁ B --- obciążenie (zasilanie z -X2, złącze sterownicze
-X5)}.

\krok{2} \textbf{Gałąź wzorcowa obciążenia.} Zbuduj pionowo, od szyny L1 do szyny N:
\ui{Trójnik} na L1 $\to$ zestyk mocy \texttt{-K11} $\to$ bezpiecznik \texttt{-F12.1} $\to$
grzałka: wyszukaj \emph{grzałka} lub \emph{rezystor} (symbol prostokąta rezystora), DT
\texttt{-E1} $\to$ \ui{Trójnik} na N. Tekst \emph{1,0} (kW) obok grzałki.

\krok{3} \textbf{Powielenie $\times$11.} Zaznacz całą gałąź ramką, \ui{Edytuj} $\to$ \ui{Powiel}
(E.6), wstaw 11 kopii w równych odstępach (kolejne trzy czwórki podpinaj Trójnikami do szyn L1,
L2, L3 --- przesuń początek gałęzi na właściwą szynę). Potem dwuklikiem popraw DT i moce:
\texttt{-K12/-F12.2/-E2/1,5} \ldots \texttt{-K22/-F12.12/-E12/3,0}, zgodnie z tabelą na
wzorze. Podpis zbiorczy pod gałęziami przepisz z PDF-a.

\krok{4} \textbf{Wentylatory (poprawka v3).} W prawym dolnym rogu części B: tekst/punkt
\texttt{X4:L} $\to$ bezpiecznik \texttt{-F5} $\to$ zestyk zwierny \texttt{-K3} $\to$ czarna
skrzynka \texttt{-K2} z tekstem \emph{I$>$} (przekaźnik prądowy --- nadzoruje, że wentylatory
biorą prąd) $\to$ Trójnik $\to$ dwa silniki: hasło \emph{silnik} (symbol M w kółku), DT
\texttt{-M3} i \texttt{-M4}, wspólny powrót do \texttt{X4:N}. Obrysuj przerywaną ramką
z podpisem \emph{wentylatory modułu B (v3: -F5, -K3)}.

\krok{5} \textbf{Zasilacz -G1.} Po lewej: trzy zaciski \texttt{X4:1..3} (L/N/PE), w linii L
bezpiecznik \texttt{-F4}, dalej czarna skrzynka \texttt{-G1} z tekstem \emph{zasilacz
24 / 5 / 3,3 V}; z prawej krawędzi skrzynki dwie strzałki tekstowe: \emph{24 V $\to$ /E2.1,
cewki, -Y1, -K4} i \emph{5 / 3,3 V $\to$ -A1, czujniki}. Tu też umieść pary punktów przerwania
\texttt{+24V} i \texttt{0V} (te z E2) --- odsyłacze przy szynach E2 zaczną wskazywać na E3.

\krok{6} \textbf{Płyta logiki -A1.} Duża czarna skrzynka \texttt{-A1} z tytułem
\emph{ESP32-S3 --- płyta logiki}. Na prawej krawędzi opisz magistrale tekstami (E.6), wiersz po
wierszu jak w PDF: SPI, I2C, UART, 1-W, ADC, DI --- wraz z listą urządzeń. Na dolnej krawędzi
trzy krótkie linie DO1/DO2/DO3 (v3: rozdzielone!) prowadzące do: cewki \texttt{-K10}, skrzynki
\texttt{-A2 WDT} i cewki \texttt{-Y1}. Cewki wstawiasz hasłem \emph{cewka}; -A2 to mała czarna
skrzynka. Nad każdą cewką tekst \emph{24 V}, pod spodem podpis \emph{DO1--DO3 (opto, low-side)}.

\krok{7} \textbf{Generator -A3.} Czarna skrzynka \texttt{-A3} z tekstem jak w PDF; obok niej
para punktu przerwania \texttt{L1\_A3} z E1 (przyłącze generatora do L1 za oknem -B1).

\krok{8} \textbf{Rząd cewek stopni.} Krótka szyna \emph{24 V} $\to$ 17 małych cewek
(\texttt{-K11}\ldots\texttt{-K22}, \texttt{-K30}\ldots\texttt{-K33}, \texttt{-K3}) $\to$ szyna
\emph{0 V}: wstaw jedną cewkę, resztę Powielem, potem popraw DT. Gdy DT cewki -K11 pokryje się
z zestykiem -K11 z kroku 2 --- pojawią się wzajemne odsyłacze (to jest cała ,,magia'' EPLAN:
cewka i styki tego samego aparatu = ten sam DT). Nad rzędem umieść skrzynkę z tekstem
\emph{-U5/-U6 ULN2803 $\leftarrow$ -U3/-U4 (I2C -A1)} i podpis zbiorczy z PDF-a.

\krok{9} \textbf{Adnotacje końcowe} (E.6) i porównanie ze wzorem strona-po-stronie.

% =====================================================================
\section{Kontrola, drukowanie, eksport}

\begin{itemize}
\item \textbf{Nawigacja po odsyłaczach:} kliknij prawym na zestyku lub cewce --- w menu
kontekstowym jest polecenie przejścia do elementu powiązanego (\szukaj{Przejdź do} pokaże
dostępne warianty). Tak sprawdzisz każdą parę cewka--styki (K1, K2, K3, K4, K10, RB, K11--K33)
i wszystkie punkty przerwania.
\item \textbf{Weryfikacja projektu:} \szukaj{Sprawdź} --- uruchom sprawdzenie projektu
i przejrzyj komunikaty (np. niesparowane punkty przerwania, zdublowane DT). Komunikaty otwierają
się dwuklikiem wprost na winnym elemencie.
\item \textbf{Eksport PDF:} karta \ui{Plik} $\to$ pozycja eksportu/publikacji PDF
(najszybciej: \szukaj{PDF}). Zaznacz wszystkie strony projektu. \emph{Uwaga:} wersja Education
dodaje do wydruków znak wodny --- to normalne i nie da się go wyłączyć.
\item \textbf{Drukowanie:} \ui{Plik} $\to$ \ui{Drukuj} (Ctrl+P).
\end{itemize}

% =====================================================================
\section{Najczęstsze problemy początkującego}

\begin{longtable}{@{}p{0.42\textwidth}p{0.54\textwidth}@{}}
\toprule
\textbf{Objaw} & \textbf{Przyczyna / rozwiązanie}\\
\midrule
\endhead
Symbole się nie łączą & przyłącza nie leżą w jednej osi siatki --- włącz przyciąganie (pasek
stanu, na dole), przesuń symbol; sprawdź, czy siatka to 4 mm\\
Nie mogę znaleźć przycisku z przewodnika & wpisz nazwę w \ui{Szukaj polecenia} (lupka na górze)
--- to zawsze działa; nazwy grup wstążki bywają skrócone przy wąskim oknie\\
Zamknąłem panel Strony / Centrum wprowadzania & \szukaj{Strony} lub \szukaj{Centrum
wprowadzania} --- panel wróci\\
Symbol wstawia się obrócony & podczas wstawiania naciskaj \textbf{Tab} (warianty A--D symbolu)\\
Lustro styków nie pokazuje się pod cewką & upewnij się, że styki mają \emph{identyczny} DT jak
cewka (łącznie z myślnikiem); potem właściwości cewki $\to$ ustawienie lustra styków\\
Punkt przerwania zgłasza brak pary & druga strona musi mieć punkt o \emph{tej samej} nazwie;
literówka = brak pary\\
Chcę cofnąć kilka operacji & Ctrl+Z wielokrotnie (pasek szybkiego dostępu, lewy górny róg)\\
Przypadkowo rysuję na złej stronie & patrz na zakładkę nad edytorem (np. \texttt{=CA1+EAA/2})
--- dwuklik właściwej strony w nawigatorze\\
\bottomrule
\end{longtable}

\section{Co dalej}
Po narysowaniu E1--E3: uzupełnij stronę tytułową (dwuklik na \texttt{/1}; pola pochodzą
z właściwości projektu --- \szukaj{Właściwości projektu}), uruchom weryfikację, wyeksportuj PDF
i porównaj z \texttt{schemat\_elektryczny\_v3\_EPLAN.pdf}. Dalsze etapy (dobór artykułów z bazy,
zestawienia aparatów, numeracja przewodów) wykraczają poza ten przewodnik --- to naturalny
następny krok po opanowaniu rysowania.

\end{document}
